\chapter{引言}

\section{研究背景和意义}

随着数字经济的快速发展，数据要素化已成为推动经济社会发展的重要驱动力。
在分布式环境下，数据的安全流通和共享变得愈发重要，
尤其是在涉及多个参与节点的场景中，如何保证数据传输和存储过程的安全性成为亟待解决的关键问题。

如何结合分布式技术完成数据的安全共享一直是热门研究方向，
自从2008年Nakamoto
\cite{nakamoto2008bitcoin}
（中本聪）发表了一篇《比特币：一种点对点的电子现金系统》后，
吸引了无数来自不同学术领域的学者的关注。

当前，传统的数据安全授权方案在分布式环境下面临诸多挑战：
首先，单一的授权节点容易成为系统安全的薄弱环节；
其次，现有的授权机制难以满足分布式环境下多方协同的需求；
最后，在确保数据安全的同时如何保持系统效率，也是一个重要问题。

本课题的研究对推动数据要素安全流通具有重要的理论和实践意义。
在理论层面，本研究将门限密码学与代理重加密技术相结合，
拓展了传统代理重加密的应用场景。在实践层面，本研究提出的基于国密算法的分布式代理重加密方案，
为解决分布式环境下的数据安全授权问题提供了可行的技术路径。

\section{国内外研究现状}

\subsection{国外研究现状}

在国外，代理重加密技术的研究始于Blaze\cite{blaze1998divertible}等人1998年提出的概念。
近年来，随着分布式系统的普及，国际学术界对分布式环境下的数据安全授权展开了深入研究。主要包括：

1）基于属性的代理重加密方案，实现了更细粒度的访问控制

2）门限代理重加密技术，提高了系统的容错性和可靠性

3）基于格密码的代理重加密方案，为抵抗量子计算攻击提供了可能

\subsection{国内研究现状}

国内在数据安全授权领域的研究主要集中在以下方面：

1）国密算法在代理重加密中的应用研究

2）基于区块链的分布式授权机制研究

3）面向特定领域的数据安全共享方案研究

然而，目前的研究仍存在一些不足，如计算效率有待提高、安全性证明不够完备等问题。

\section{主要内容和工作安排}

本文共分为五章，具体安排如下：

第1章为引言，介绍研究背景、意义及国内外研究现状。

第2章为理论基础，阐述代理重加密、门限密码和国密算法的基本原理。

第3章为系统设计，详细描述分布式环境下基于国密的代理重加密方案。

第4章为实验验证，通过多组对比实验验证系统的性能和安全性。

第5章为总结与展望，总结研究成果并提出未来的研究方向。

本研究预期在保证数据安全性的同时，实现高效的分布式数据授权管理，为解决数据要素化背景下的数据安全流通问题提供新的技术方案。
